\tzpanchor{0368}
\tzpentryheader{0368}{Extended Helicity Functional}

\begingroup\small
\begingroup\scriptsize
\setlength{\tabcolsep}{4pt}
\renewcommand{\arraystretch}{0.92}
\noindent\begin{tabularx}{\linewidth}{@{}p{0.24\linewidth}p{0.36\linewidth}X@{}}
\textbf{UUID} & \multicolumn{2}{l@{}}{\tzpcode{5527b462-854e-57a8-a437-dc8984272b52}}\\
\textbf{PiID} & \tzpcode{0548820466521} & \textbf{TZPi}\quad \tzpcode{0}\quad \textbf{PiPE}\quad \tzpcode{375}\\
\textbf{RTEID} & \textbf{Poloidal $\theta$}\quad \tzpcode{1429.0257 rad} & \textbf{Toroidal $\phi$}\quad \tzpcode{02.3546 rad}\\
\textbf{TZPID} & \tzpcode{ID0368} & \textbf{Node Dependencies}\quad \tzpidref{0367}\\
\textbf{Registry} & \tzpcode{registry\_id\_0368} & \textbf{Scale}\quad transfer\\
\textbf{LOC Classification} & QA641 manifolds and topology & QC174.17 mathematical physics\\
\textbf{arXiv Classification} & Primary: math-ph & Cross-list: gr-qc, math.DG\\
\textbf{Primary Support} & Variation Law & Support Relation\\
\textbf{Closure Support} & Closure Relation & \\
\end{tabularx}\par
\endgroup
\endgroup

\tzpsection{Canonical Statement}
Helicity measures the linkage between magnetic field lines, corrected for rotation and contributions from extra dimensions.

\tzpsection{Canonical Equation}
\[
\mathcal{H}_{\mathrm{ext}}=\int_{\Omega}\mathbf{A}\cdot\mathbf{B}\,d^{3}x+\int_{\partial\Omega}\alpha\,dS
\]
\[
\mathbf{B}=\nabla\times\mathbf{A}
\]

\begin{minipage}[t]{0.49\linewidth}
\tzpsection{Canonical Symbol Key}
\begingroup
\small
\raggedright
\textbf{Source equation symbols.}\par
\begin{itemize}[leftmargin=1.35em,itemsep=1pt,topsep=1pt,parsep=0pt]
    \item $\mathcal{H}_{\mathrm{ext}}$: extended helicity functional.
    \item $\Omega$: domain, angular, or boundary operator associated with the equation.
    \item $A$: vector potential whose curl defines the magnetic field.
    \item $B$: magnetic field defined by the curl of the vector potential.
\end{itemize}
\endgroup
\end{minipage}\hfill
\begin{minipage}[t]{0.47\linewidth}
\tzpsection{Wolfram Form}
\begingroup
\scriptsize
\raggedright
\textbf{Source Wolfram model.}\par
\begin{Verbatim}[fontsize=\tiny,breaklines=true,breakanywhere=true]
(* TZPID ID0368 generated source Wolfram model *)
ClearAll[K0368, Cr0368, T, R, A, W, I, N];
eq0368$1 = HoldForm[HExt == IntegralOmegaA*B*d^3*x + IntegralPartialOmegaalpha*dS];
eq0368$2 = HoldForm[B == Grad*A];
eq0368$3 = HoldForm[(d/dt)*HExt == 0];

K0368 = HoldForm[{eq0368$1, eq0368$2, eq0368$3}];

N = "normalized TRAWIN closure object for Extended Helicity Functional";
I = "Closure Relation integration/comparison layer";
W = "Support Relation wave or operator carrier";
A = "Variation Law admissible action/amplitude condition";
R = "Support Relation rotation/curl preservation layer";
T = "Variation Law local source condition";

Cr0368[expr_] := HoldForm[N[I[W[A[R[T[expr]]]]]]];

Result0368 = Cr0368[K0368];
\end{Verbatim}
\endgroup
\end{minipage}

\par\vspace{0.45\baselineskip}
\noindent\begin{minipage}[t]{\linewidth}
\begingroup
\small
\raggedright
\textbf{TRAWIN Operator Key.}\par
\noindent\textbf{Time/localization operator:} $\mathsf{T}\!\left[\mathrm{local}\right]$ Variation Law local source condition.\par
\noindent\textbf{Rotation/curl operator:} $\mathsf{R}\!\left[\nabla\times\right]$ Support Relation rotation/curl preservation layer.\par
\noindent\textbf{Action/amplitude operator:} $\mathsf{A}\!\left[\mathrm{admissible}\right]$ Variation Law admissible action/amplitude condition.\par
\noindent\textbf{Wave/d'Alembertian operator:} $\mathsf{W}\!\left[\Box\right]$ Support Relation wave or operator carrier.\par
\noindent\textbf{Integration operator:} $\mathsf{I}\!\left[\int\right]$ Closure Relation integration/comparison layer.\par
\noindent\textbf{Normalization operator:} $\mathsf{N}\!\left[\mathrm{normalized}\right]$ normalized TRAWIN closure object for Extended Helicity Functional.\par
\endgroup
\end{minipage}

\clearpage
\noindent\textbf{[ID 0368] Extended Helicity Functional}\par
\vspace{0.35em}
\tzpsection{Derivation Layer}
\paragraph{Primary Equation.}
\[
\mathbf{B}=\nabla\times\mathbf{A}
\]
\[
\mathbf{E}=-\dot{\mathbf{A}}-\nabla\Phi
\]

\paragraph{Primary Derivation.}
The recovered conductive-media passage supplies the potential-to-field relations used as support for the helicity functional: the magnetic field is the curl of the vector potential, while the electric field is generated by its time derivative together with the scalar-potential gradient.

\paragraph{Interpretation.}
These relations ground the field variables entering extended helicity. They do not independently prove conservation of the full boundary-corrected helicity functional.

\par\vspace{0.35em}
\noindent\textbf{Derivable Consequences.}\quad
\begingroup
\footnotesize
\raggedright
The potential-to-field relations anchor magnetic structure before the progression applies them within the extended helicity functional.
\par\endgroup

\tzpsection{Equation Support}
\noindent\textbf{1. Variation Law}\par
\[
\mathcal{H}_{\mathrm{ext}}=\int_{\Omega}\mathbf{A}\cdot\mathbf{B}\,d^{3}x+\int_{\partial\Omega}\alpha\,dS
\]
\noindent This fixes the first explicit relation carried by the entry rather than a generic certificate record normalization.\par
\noindent\textbf{2. Support Relation}\par
\[
\mathbf{B}=\nabla\times\mathbf{A}
\]
\noindent This adds the next operational equation that advances the entry beyond its opening definition.\par
\noindent\textbf{3. Closure Relation}\par
\[
\frac{d}{dt}\mathcal{H}_{\mathrm{ext}}=0
\]
\noindent This closes the progression with an actual admissibility or closure relation instead of recycling the normalization shell.\par

\clearpage
\noindent\textbf{[ID 0368] Extended Helicity Functional}\par
\vspace{0.35em}
\tzpsection{Dictionary}
\begingroup\small

\tzpsection{Definition}

The extended helicity functional is grounded by potential-to-field relations: a vector potential produces a magnetic field through its curl and participates in the associated electric-field relation.

\tzpsection{Contextual Meaning}

Helicity measures organized linkage in a field configuration. The magnetic field in that reading must be tied to a generating potential, so \(\mathbf{B}=\nabla\times\mathbf{A}\) establishes the field structure before boundary extensions are interpreted.

\tzpsection{Formal Statement}

\[
\mathbf{B}=\nabla\times\mathbf{A}, \qquad \mathbf{E}=-\dot{\mathbf{A}}-\nabla\Phi
\]

\tzpsection{Derivable Consequences}

Once \(\mathbf{A}\) is specified, \(\mathbf{B}\) is fixed by its curl. Adding a gradient leaves that curl unchanged under the usual gauge reading; complete helicity interpretation remains conditional.

\tzpsection{Lemma Support}

The approved passage supports potential-to-field relations in conductive media and defines \(\mathbf{A}\) and \(\mathbf{B}\). It does not certify every extended-helicity claim.

\tzpsection{Technical Interpretation}

Curl extracts rotational field content from the potential; the potential-field pairing is the mathematical carrier for linkage-sensitive field structure.

\tzpsection{Equation Grounding}

\(\mathbf{A}\) is the vector potential whose curl defines the magnetic field; \(\mathbf{B}\) is the magnetic field defined by that curl. The associated \(\mathbf{E}\) and \(\Phi\) terms are displayed for relation context but await independent symbol-level approval.

\tzpsection{Interpretive Role}

This entry supplies the field-definition bridge required before later helicity and linkage interpretations are stated.

\tzpsection{TZP Thesaurus}

\noindent \textbf{Vector Potential}: A generating field whose curl supplies the magnetic field. \enspace\textbullet\enspace \textbf{Magnetic Curl Relation}: The defining relation that connects potential structure to magnetic structure. \enspace\textbullet\enspace \textbf{Helicity Carrier}: The potential-field pairing used to represent linkage-sensitive organization.\par

\endgroup

\tzpsection{Encyclopedia Mathematica}
\ifTZPDarkEdition
\tzpdictionaryvisualband{D:/TZPID/kdp_workflow/build/tikz_figures/ID0368_dictionary_figure_dark.pdf}
\fi

\clearpage
\begingroup\tzpsupportsetup
\noindent\textbf{\fontsize{10.8pt}{12.8pt}\selectfont [ID 0368] Constructive Logic and Proof Certificate}\par
\noindent\textbf{\fontsize{10pt}{12pt}\selectfont Extended Helicity Functional}\par
\vspace{0.45em}
\begingroup
\footnotesize
\setlength{\parskip}{0.22em}
\setlength{\parindent}{0pt}
\noindent\textbf{Curry--Howard--Lambek Interpretation}\par
Under the Curry--Howard--Lambek correspondence, this registry entry is read as a typed proof
object: the stated source condition supplies the proposition, the derivation supplies the
morphism, and the proof certificate records the machine handles needed to check the obligation
in the formal registry.

\vspace{0.45em}
\noindent\textbf{CHL Proof}\par
\textbf{Statement:} helicity measures the linkage between magnetic field lines, corrected for rotation and
contributions from extra dimensions

\textbf{Logic Validation:}\\
\textbf{Proposition:} ``Extended Helicity Functional is valid as a formal dynamical law when the Hamiltonian or
energy operator is well defined on the stated state space.''\\
\textbf{Proof:} The source statement identifies an energy, Hamiltonian, or vacuum-sector mechanism. The
CHL proof reads it as an operator claim: if the state space and localized terms are
defined, then the evolution or extraction channel is formally meaningful.

\textbf{VERDICT:} \ensuremath{\checkmark}\ \textbf{TRUE} --- formal dynamical-operator validation

\textbf{Formal Status:} \ensuremath{\checkmark}\ THEORETICAL -- source-backed CHL proof obligation

\textbf{Physical Status:} THEORETICAL -- requires model-specific empirical interpretation
\par\vspace{0.45em}
\noindent{\tiny\textbf{Isabelle/HOL.} \tzpplaincode{physics\_validity\_from\_model, external\_validity\_package, registry\_number\_matches, equation\_count\_matches}}\par
\ifTZPDarkEdition
\tzpproofvisualband{D:/TZPID/kdp_workflow/build/tikz_figures/ID0368_proof_certificate_figure_dark.pdf}
\fi
\endgroup
\endgroup
